\section{Paying for an increase of potential}\label{sec:potential-step}

We prove Proposition~\ref{prop:potential-step}. For a transition of
width one, two bounds on the orders of the endpoint groups restrict the
possible layer changes. A scalar optimization then shows that the
certificate cost pays for the increase of potential. For width two,
the gate is paid for directly by the admissible-column estimate.

\subsection{Layer changes and order budgets}

Fix a transition of width one from a nonsingular matrix \(C\) to a
nonsingular matrix \(B\), and let \(Q\) be the common quotient obtained
from \(G_C\) by the new column and from \(G_B\) by the new coordinate.
For an old layer \((p,e)\), write
\[
 k=k_C(p,e),\qquad k'=k_Q(p,e).
\]
The interlacing of layer heights gives
\[
 k'\in\{k-1,k\},\qquad k_B(p,e)\in\{k',k'+1\}.
\]
If \(k_C(p,e)=0\), the same statement gives \(k_B(p,e)\le1\), so this
layer contributes nothing to either potential.

Use the weights \(\nu_p\) in all sums below. In the height-one category,
mark exactly the layers that reach height two; in the category of old
height at least two, mark every retained layer. Let \(g,v_0\) be the
marked and unmarked weights in the first category, and \(a,u\) those in
the second. The following six types define \(L\) and \(P\); the last
column bounds the potential change per unit weight.
\begin{center}
\begin{tabular}{@{}ccccc@{}}
\toprule
weight & \shortstack{old\\height} & retention status
 & \shortstack{destination\\height}
 & \shortstack{potential change\\per unit, at most}\\
\midrule
\(g\) & \(1\) & retained & \(2\) & \(\beta\)\\
\(v_0\) & \(1\) & retention may occur & \(\le1\) & \(0\)\\
\(a-L\) & \(k\ge2\) & retained & \(k+1\) & \(6-\beta\)\\
\(L\) & \(k\ge2\) & retained & \(k\) & \(0\)\\
\(u-P\) & \(k\ge2\) & not retained & \(k-1\) & \(-\beta\)\\
\(P\) & \(k\ge2\) & not retained & \(k\) & \(0\)\\
\bottomrule
\end{tabular}
\end{center}
The \(g\)-row is retained because \(k_B(p,e)=2\) forces \(k'=1\).
The potential bounds follow from \(\phi(2)=\beta=7/5\) and
\(\phi(k)=2k\) for \(k\ge3\): an increase from height at least two
costs at most \(6-\beta\) per unit weight, and a decrease saves at
least \(\beta\). Layers that stay at the same height make no change.

Both categories consist of whole plateaus determined by \(G_C\). The
marks are prefixes on these plateaus: this is the retention property
for the second category, and for the first it follows from the fact
that \(k_B(p,e)\) is nonincreasing in \(e\). Thus these marks have the
form required in Proposition~\ref{prop:certificate-pricing}.

We claim
\begin{equation}\label{eq:flat-compensation}
 P\le L,\qquad a_0:=a-L+P\in[0,a].
\end{equation}
Here \(a_0\) is an effective scalar mass; the actual marked family in
the second category has mass \(a\). Thus every \(0\le y\le a_0\le a\) is a
valid target for that marked family.
To prove the first assertion, fix \(p\) and let \(t\) be the last
exponent at which \(k_C(p,e)\ge2\), taking \(t=0\) if there is none.
Attach a subscript \(p\) to the contribution of this prime to each
weight sum. In the notation of Lemma~\ref{lem:layer-inequality}, a
retained layer in this range increases the height by one unless it is
flat; an unmarked layer decreases it by one unless it persists.
Consequently that Lemma, with \(h=1\), gives
\[
 (a_p-L_p)-(u_p-P_p)
 =\nu_p\bigl(F_B(p,t)-F_C(p,t)\bigr)
 \le \nu_p(t-2U_t)=a_p-u_p.
\]
The last equality holds because \(U_t\) counts the unmarked layers
in this range. Hence \(P_p\le L_p\), and summing over primes proves
\(P\le L\). Also \(L\le a\), so \(a_0=a-L+P\) is nonnegative and at
most \(a\).

Summing the potential bounds in the table gives
\begin{equation}\label{eq:potential-increment}
 \begin{aligned}
 \Psi(B)-\Psi(C)
 &\le \beta g+(6-\beta)(a-L)-\beta(u-P)\\
 &\le I:=\beta g+(6-\beta)a_0-\beta u.
 \end{aligned}
\end{equation}
Indeed, the first line equals \(I-(6-2\beta)P\), and \(6-2\beta>0\).

The total height times weight is at most \(1/2\) at both endpoints.
At the \(B\) endpoint it is at least
\[
 2g+3(a-L)+2L+(u-P)+2P
 =2g+3a+u-L+P
 \ge 2g+3a_0+u.
\]
At the \(C\) endpoint it is at least
\(g+v_0+2a+2u\), which is at least \(g+v_0+2a_0+2u\).
We have therefore obtained the two budgets
\begin{equation}\label{eq:potential-budgets}
 2g+3a_0+u\le\frac12,\qquad
 g+v_0+2a_0+2u\le\frac12.
\end{equation}

\subsection{Scalar bound and conditional potential estimate}

\begin{lemma}[Budget optimization]\label{lem:budget-optimization}
Let \(c\) and \(\mathcal E\) be defined by
Equations~\eqref{eq:cube-profile} and \eqref{eq:entropy-cost}, and let
\(\beta=7/5\). If \(g,v_0,a_0,u\ge0\) satisfy
\[
 2g+3a_0+u\le\frac12,\qquad
 g+v_0+2a_0+2u\le\frac12,
\]
then there are \(x\in[0,g]\) and \(y\in[0,a_0]\) for which
\begin{equation}\label{eq:budget-primal}
 c(x+2y)-\mathcal E(g,v_0,x)-\mathcal E(a_0,u,y)
 \ge \beta g+(6-\beta)a_0-\beta u.
\end{equation}
\end{lemma}

The inequality says that certificate concentration, after paying both
entropy costs, covers the upper bound \(I\) for the potential increase.
The proof of the lemma follows in the next subsection.

\begin{proof}[Proof of Proposition~\ref{prop:potential-step}]
Fix the desired precision \(A>0\), and set \(A_1=A+2\). For width
one, take a nonsingular starting matrix \(C\) in the typical set of
Proposition~\ref{prop:certificate-pricing} at precision \(A_1\) for
\((w_1,w_2)=(1,2)\). Extend the marking rule above to singular
outcomes by leaving all layers unmarked. Apply that Proposition with
exposed information \(C\), starting quotient \(G_C\), and the subevent
\(\{\det B\ne0\}\).
The tested column is uniform conditional on \(C\), and the marks
defined above are allowed to depend also on the new diagonal bit.
The Proposition gives at most
\(2^{N\ell^{-A_1}}\) classes. In each nonempty class
\(\mathcal C\), choose an outcome \(B_{\mathcal C}\) maximizing
\(\Psi(B)\) in that class. Such a representative exists because
the outcome space is finite.

For this representative, form \(g,v_0,a,u,L,P,a_0\) and \(I\) as
above. Lemma~\ref{lem:budget-optimization} supplies \(x\in[0,g]\)
and \(y\in[0,a_0]\). These are admissible certificate target weights
for the categories of minimum heights one and two, since \(a_0\le a\).
The monotonicity of \(\mathcal E\) in its first argument gives
\[
 \begin{aligned}
 c(x+2y)-\mathcal E(g,v_0,x)-\mathcal E(a,u,y)
 &\ge c(x+2y)-\mathcal E(g,v_0,x)-\mathcal E(a_0,u,y)\\
 &\ge I.
 \end{aligned}
\]
Proposition~\ref{prop:certificate-pricing}, applied to this
representative in the restricted outcome space, therefore yields
\[
 \Pr(\mathcal C\mid C)\le2^{-N(I-\ell^{-A_1})}.
\]
Since \(\Psi(B_{\mathcal C})-\Psi(C)\le I\) by
Equation~\eqref{eq:potential-increment}, the contribution of this
class satisfies
\[
 \E\!\left[\ind_{\mathcal C}\,2^{N\Psi(B)}\mid C\right]
 \le 2^{N\Psi(B_{\mathcal C})}\Pr(\mathcal C\mid C)
 \le 2^{N(\Psi(C)+\ell^{-A_1})}.
\]
Summing over the classes proves
\[
 \E\!\left[\ind_{\{\det B\ne0\}}\,2^{N\Psi(B)}\mid C\right]
 \le2^{N(\Psi(C)+2\ell^{-A_1})}
 \le2^{N(\Psi(C)+\ell^{-A})}
\]
for all sufficiently large \(N\).

For width two, take \(C\) also in the typical set of
Proposition~\ref{prop:admissible-potential} at precision \(A_1\),
and let \(m\) be its size. Conditional on \(C\), the two cross
columns \(v_1,v_2\) are independent and uniform on \(\bits^m\).
For the gate \(\mathcal G_2=\{v_1,v_2\in\mathcal S(C)\}\), that
Proposition gives
\[
 \Pr(\mathcal G_2\mid C)
 =2^{-2m}|\mathcal S(C)|^2
 \le2^{-2m+2N(\Psi(C)+\ell^{-A_1})}.
\]
On nonsingular outcomes \(\Psi(B)\le1\). Consequently
\[
 \begin{aligned}
 \E\!\left[\ind_{\{\det B\ne0\}}\ind_{\mathcal G_2}
                2^{N\Psi(B)}\mid C\right]
 &\le2^{N-2m+2N(\Psi(C)+\ell^{-A_1})}\\
 &\le2^{N(\Psi(C)+2(1-m/N)+2\ell^{-A_1})}\\
 &\le2^{N(\Psi(C)+\ell^{-A})}.
 \end{aligned}
\]
The second line uses \(\Psi(C)\le1\), and the last uses
\(1-m/N\le N^{-3/10}=o(\ell^{-A})\) together with \(A_1=A+2\).
Intersecting the two typical sets proves both conditional assertions
of Proposition~\ref{prop:potential-step}, with arbitrary fixed
precision and with the stated nonsingularity and gate indicators.
\end{proof}

\subsection{Proof of the budget optimization}

\begin{proof}[Proof of Lemma~\ref{lem:budget-optimization}]
The second budget ensures \(x+2y\le1/2\) throughout the rectangle
under consideration. Sections~\ref{sec:cube} and
\ref{sec:certificate-pricing} show that \(c\) is nondecreasing and
concave, while \(\mathcal E\) is continuous, nonnegative, convex in its
target, homogeneous of degree one in all size variables, and
nonincreasing in its first argument for fixed \(u,x\).
Equation~\eqref{eq:cube-profile} further shows that \(c\) is
differentiable on \((0,1/2)\), including the join at \(q=1/4\), has
left derivative \(1\) at \(1/2\), and satisfies
\(c(q)/q\to\infty\) as \(q\downarrow0\).

\medskip\noindent\textbf{Supporting slope.}
For \(w,t\ge0\), put
\[
 c^*(w)=\max_{0\le q\le1/2}\bigl(c(q)-wq\bigr),\qquad
 \mathcal F_t(s,u)=\max_{0\le x\le s}
       \bigl(tx-\mathcal E(s,u,x)\bigr).
\]
Write \(I=\beta g+(6-\beta)a_0-\beta u\). We will prove that, for
every \(w\ge0\),
\begin{equation}\label{eq:budget-dual}
 c^*(w)+\mathcal F_w(g,v_0)+\mathcal F_{2w}(a_0,u)\ge I.
\end{equation}
Here is a justification that this suffices. If \(g=a_0=0\), taking
\(x=y=0\) proves Equation~\eqref{eq:budget-primal} immediately.
Otherwise let \((x_*,y_*)\) maximize its left side on the compact
rectangle, and put \(q_*=x_*+2y_*\). At zero the objective is zero.
For a positive one of \(g,a_0\), its entropy cost is \(O(x)\) or
\(O(y)\) near zero, by the derivative of \(H\) recorded in
Section~\ref{sec:certificate-pricing}; the fact that
\(c(q)/q\to\infty\) then shows that the maximum is positive.
Thus \(q_*>0\).

Take \(w=c'(q_*)\) if \(q_*<1/2\), and the left derivative if
\(q_*=1/2\). This is a nonnegative supporting slope, so
\(c^*(w)=c(q_*)-wq_*\). Let
\[
 K(x,y)=\mathcal E(g,v_0,x)+\mathcal E(a_0,u,y).
\]
We claim that \((x_*,y_*)\) maximizes \(w(x+2y)-K(x,y)\) on the
same rectangle. If some \((x,y)\) gave a strictly larger value, take
the point a fraction \(\theta\) along the segment from
\((x_*,y_*)\) to \((x,y)\). Convexity of \(K\) and the one-sided
differentiability of \(c\) in this feasible direction give, as
\(\theta\downarrow0\), an increase in the original objective of at
least
\[
 \theta\bigl(w(x+2y-q_*)-K(x,y)+K(x_*,y_*)\bigr)+o(\theta)>0,
\]
a contradiction. The claimed maximum separates into its two
coordinates. Hence the maximum of the left side of
Equation~\eqref{eq:budget-primal} is exactly
\(c^*(w)+\mathcal F_w(g,v_0)+\mathcal F_{2w}(a_0,u)\) at this
supporting slope. Equation~\eqref{eq:budget-dual} proves the desired
assertion.

\medskip\noindent\textbf{Entropy minimum and criterion.}
We use the function \(H\) from Equation~\eqref{eq:entropy-difference}
to compute a lower bound for the two \(\mathcal F\) terms.
For \(t>0\) set
\[
 D(t)=-\log_2(1-2^{-t}).
\]
For \(r\ge0\), maximizing \(H(z,r)-tz\) over \(0\le z\le1\)
at \(z=\min\{1,r/(2^t-1)\}\) gives
\[
 \mathcal F_t(1,r)=
 \begin{cases}
 t-H(1,r)+D(t)r,&0\le r\le2^t-1,\\
 0,&r\ge2^t-1.
 \end{cases}
\]
For \(h\ge0\), the derivative on the first branch after adding \(hr\)
is \(D(t)+h-\log_2(1+1/r)\). It is increasing in \(r\) and vanishes
at \(r=(2^{D(t)+h}-1)^{-1}\le2^t-1\), with equality when \(h=0\).
On the second branch \(hr\) is nondecreasing. Evaluation at this
minimum therefore gives
\begin{equation}\label{eq:entropy-dual-minimum}
 \begin{aligned}
 \min_{r\ge0}\bigl(\mathcal F_t(1,r)+hr\bigr)
 &=t-D(D(t)+h)\\
 &=\mathcal L(t,h):=
   \log_2\!\left(1+(2^t-1)(1-2^{-h})\right).
 \end{aligned}
\end{equation}
For \(t=0\), both the minimum and \(\mathcal L(0,h)\) are zero,
so the equality between the first and last expressions holds there
as well; the middle expression is used only for \(t>0\).

For a given \(w\), it is enough to find \(\lambda,\mu\ge0\) such
that
\begin{equation}\label{eq:budget-multipliers}
 \begin{aligned}
 2\lambda+\mu+\mathcal L(w,\mu)&\ge\beta,\\
 3\lambda+2\mu+
   \mathcal L(2w,\beta+\lambda+2\mu)&\ge6-\beta,\\
 (\lambda+\mu)/2&\le c^*(w).
 \end{aligned}
\end{equation}
Indeed, homogeneity and Equation~\eqref{eq:entropy-dual-minimum}
imply
\[
 \begin{aligned}
 \mathcal F_w(g,v_0)+\mu v_0&\ge g\,\mathcal L(w,\mu),\\
 \mathcal F_{2w}(a_0,u)+(\beta+\lambda+2\mu)u
   &\ge a_0\,\mathcal L(2w,\beta+\lambda+2\mu).
 \end{aligned}
\]
These remain true when \(g=0\) or \(a_0=0\). Using the first two
inequalities in Equation~\eqref{eq:budget-multipliers} now gives
\[
 \mathcal F_w(g,v_0)+\mathcal F_{2w}(a_0,u)
 +\lambda(2g+3a_0+u)+\mu(g+v_0+2a_0+2u)\ge I.
\]
The two budgets and the third inequality in
Equation~\eqref{eq:budget-multipliers} yield
Equation~\eqref{eq:budget-dual}.

\medskip\noindent\textbf{Exact verification.}
We verify Equation~\eqref{eq:budget-multipliers} with three explicit
choices. The function \(D\) is decreasing, and, for \(h>0\),
\begin{equation}\label{eq:L-linear-bound}
 \mathcal L(t,h)\ge t-D(h).
\end{equation}
This follows by dropping the positive term \(2^{-h}\) from
\(2^{\mathcal L(t,h)}=2^t(1-2^{-h})+2^{-h}\).
Only one numerical estimate for \(D\) is needed:
\begin{equation}\label{eq:D-uniform-bound}
 D(21/10)<2/5.
\end{equation}
Indeed, if \(v=2^{-1/10}\), then \(v<14/15\) because
\(15^{10}<2\cdot14^{10}\). Thus
\[
 2^{-21/10}+2^{-2/5}
 =\frac v4+v^4
 <\frac7{30}+\left(\frac{14}{15}\right)^4
 =\frac{100457}{101250}<1,
\]
which is equivalent to Equation~\eqref{eq:D-uniform-bound}.

We will use the support bounds
\begin{equation}\label{eq:dual-support-lines}
 2c^*(w)\ge
 \max\left\{2-w,\ \frac32-\frac w2,\
             \frac{23}{24}-\frac w6,\ \frac2{w+2}\right\}.
\end{equation}
The first three follow by testing \(q=1/2,1/4,1/12\) in the
definition of \(c^*\). The first two values are \(c(1/2)=1\) and
\(c(1/4)=3/4\). At \(q=1/12\), the argument of \(z\) is \(1/5\),
the square root in its definition is \(3/5\), and
\(c(1/12)=1-(5/6)(5/8)=23/48\).
For the last bound, when \(0\le q\le1/4\), put
\[
 r=\frac{2\sqrt{q(1-3q)}}{1-2q}.
\]
Then
\[
 r^2-4q=\frac{4q^2(1-4q)}{(1-2q)^2}\ge0.
\]
The expression \((r+2q)/(1+r)\) increases in \(r\), so
\[
 c(q)=\frac{r+2q}{1+r}
 \ge\frac{2\sqrt q+2q}{1+2\sqrt q}
 \ge2\sqrt q-2q.
\]
The difference in the last inequality is
\(4q\sqrt q/(1+2\sqrt q)\). Testing
\(q=(w+2)^{-2}\le1/4\) now gives
\[
 2c^*(w)\ge2\bigl(c(q)-wq\bigr)
 \ge2\left(\frac2{w+2}-\frac{w+2}{(w+2)^2}\right)
 =\frac2{w+2},
\]
including at \(w=0\).

One global estimate handles the middle range:
\begin{equation}\label{eq:L-quarter-bound}
 \mathcal L(w,1/4)>\frac{w-1}{2}\qquad(w\ge0).
\end{equation}
To prove it, let \(b=2^{-1/4}\). The integer inequality
\(20^4<2\cdot17^4\) gives \(1/2<b<17/20\). Since
\(b(1-b)\) decreases for \(b\ge1/2\),
\[
 b(1-b)>\frac{17}{20}\frac3{20}
 =\frac{51}{400}>\frac18.
\]
Consequently the arithmetic--geometric mean inequality gives
\[
 2^{\mathcal L(w,1/4)}=b+(1-b)2^w
 \ge2\sqrt{b(1-b)2^w}>2^{(w-1)/2},
\]
which proves Equation~\eqref{eq:L-quarter-bound}.

We use the following choices, writing \(h=\beta+\lambda+2\mu\):
\[
 (\lambda,\mu)=
 \begin{cases}
 \displaystyle\left(\max\left\{2-w,\frac32-\frac w2\right\},0\right),
     &\displaystyle 0\le w\le\frac85,\\[6pt]
 \displaystyle\left(\frac{17}{24}-\frac w6,\frac14\right),
     &\displaystyle \frac85\le w\le3,\\[6pt]
 (0,3\,2^{-w}),&w\ge3.
 \end{cases}
\]
Either choice works at a shared endpoint.

For \(0\le w\le8/5\), put
\(S=\max\{2-w,3/2-w/2\}\), so \(\lambda=S\) and \(\mu=0\).
Then \(S\ge7/10\), so \(h\ge21/10\) and
\[
 2\lambda+\mathcal L(w,0)=2S\ge7/5.
\]
By Equations~\eqref{eq:L-linear-bound} and
\eqref{eq:D-uniform-bound}, the second multiplier expression exceeds
\[
 3S+2w-\frac25
 =\max\left\{\frac{28}{5}-w,\frac{41}{10}+\frac w2\right\}
 \ge\frac{23}{5}.
\]
For the last inequality, use the first term when \(w\le1\) and the
second when \(w\ge1\). The choices are nonnegative, and the third
multiplier inequality, in the equivalent form
\(\lambda+\mu\le2c^*(w)\), follows directly from
Equation~\eqref{eq:dual-support-lines}.

For \(8/5\le w\le3\), take
\(\lambda=17/24-w/6\) and \(\mu=1/4\). These are nonnegative, and
\[
 h\ge\frac75+\frac5{24}+\frac12
 =\frac{253}{120}>\frac{21}{10}.
\]
By Equation~\eqref{eq:L-quarter-bound}, the first multiplier
expression exceeds
\[
 2\lambda+\mu+\frac{w-1}{2}
 =\frac76+\frac w6
 \ge\frac{43}{30}>\frac75.
\]
By Equations~\eqref{eq:L-linear-bound} and
\eqref{eq:D-uniform-bound}, the second expression exceeds
\[
 3\lambda+2\mu+2w-\frac25
 =\frac{89}{40}+\frac32w
 \ge\frac{37}{8}>\frac{23}{5}.
\]
Finally, \(\lambda+\mu=23/24-w/6\), so
Equation~\eqref{eq:dual-support-lines} gives the third inequality.

For \(w\ge3\), take \(\lambda=0\) and \(\mu=3\,2^{-w}\), which
are nonnegative. The elementary bound
\[
 \log2=2\int_0^{1/3}\frac{dt}{1-t^2}>\frac23
\]
and \(2^\mu-1\ge\mu\log2\) give
\[
 2^{\mu+\mathcal L(w,\mu)}
 =1+2^w(2^\mu-1)
 \ge1+3\log2>3>2^{7/5}.
\]
The last inequality follows from \(3^5>2^7\). This proves the first
multiplier inequality. Here \(h\ge\beta>1\), so
\(D(h)<D(1)=1\). Equation~\eqref{eq:L-linear-bound} therefore
shows that the second expression is greater than
\[
 2w-1\ge5>\frac{23}{5}.
\]
For the third inequality, \(3\,2^{-3}=3/8<2/5\), and
\(2^w/(w+2)\) increases for \(w\ge3\), since its logarithmic
derivative is
\[
 \log2-\frac1{w+2}>\frac23-\frac15>0.
\]
It follows that \(\mu\le2/(w+2)\le2c^*(w)\) by
Equation~\eqref{eq:dual-support-lines}. This verifies
Equation~\eqref{eq:budget-multipliers} for every \(w\ge0\), and
hence proves Equation~\eqref{eq:budget-primal}.
\end{proof}
